% !TEX root = ../../main.tex
% C6.1 - Choice of technologies and models (translated). YOLO11 figures from docs.ultralytics.com/models/yolo11.
\section{Evaluation and Choice of Technologies and Models}
\label{sec:6.1}

The options were evaluated against four common criteria: meeting the requirements of Chapter~\ref{chapter:phantich}; running on a CPU-only machine with 16~GB of memory; being open-source software with a licence suited to a non-commercial project; and the developer's familiarity with them, to reduce risk within limited time. Table~\ref{tab:tech-summary} summarises the choices, and the subsections below explain the important ones.

\begin{table}[H]
\centering
\caption{Summary of technology and model choices}
\label{tab:tech-summary}
\begin{tabularx}{\textwidth}{|>{\raggedright\arraybackslash}p{3.0cm}|>{\raggedright\arraybackslash}p{3.1cm}|>{\raggedright\arraybackslash}p{3.4cm}|L|}
\hline
\textbf{Component} & \textbf{Choice} & \textbf{Alternatives considered} & \textbf{Main reason} \\
\hline
Front end & React 19, Vite, Tailwind CSS & Server-side rendering; Angular, Vue & Many interactions on one screen; front end and server separated by an API \\
\hline
Server & Python, Flask 3.1, SQLAlchemy, Alembic & Django, FastAPI & Same language as the AI models; small and explicit \\
\hline
Queue & A table in PostgreSQL & Celery with Redis or RabbitMQ & One sequential worker; few infrastructure components \\
\hline
Relational database & PostgreSQL 17 & MySQL & CHECK constraints, partial indexes, JSONB, \texttt{SKIP LOCKED} \\
\hline
Vector storage & Milvus 2.6 & pgvector, FAISS & Pre-filtering, durable storage, one collection per version \\
\hline
Frame storage & MinIO & Files on disk & Bucket-level permissions, hashes, shared with Milvus \\
\hline
Person detection & YOLO11n & YOLO11s, YOLO11m, YOLOX, Faster R-CNN & Fastest on CPU; YOLOX as a licence fallback \\
\hline
Tracking & ByteTrack (default), BoT-SORT & DeepSORT, StrongSORT & No appearance features per person per frame \\
\hline
Encoder & RaSa & Two separate models; CLIP & One shared space for the three forms of query \\
\hline
\end{tabularx}
\end{table}

\subsection{Front-End Technologies}
\label{sec:6.1.1}

The interface has many interactions on one screen, such as dragging and pasting a query image, sorting and filtering results without a new query, and opening dialogs over the list. These interactions suit a \emph{single-page application}, in which the browser loads the application once and then only exchanges data through the API, better than server-side rendering. React was chosen for its large ecosystem with ready-made libraries for routing, API calls and styling, together with Vite for development and bundling, and Tailwind CSS to keep the layout consistent and responsive on phones.

\subsection{Server-Side Technologies}
\label{sec:6.1.2}

The server and the background worker are written in Python because all the project's AI models are provided for Python. Using one language lets the server load the query encoder directly and share the data access code with the background worker. Among the three common frameworks, Django is heavier than an API-only server needs, and the asynchronous advantage of FastAPI matters little because the slowest task, video analysis, has already been moved out of the server (Section~\ref{sec:5.1.3}). Flask is small and explicit, and the developer had already used it. The queue lives in PostgreSQL instead of Celery with Redis or RabbitMQ, because there is only one sequential worker, while each additional infrastructure component costs memory on the target machine.

\subsection{Database and Storage}
\label{sec:6.1.3}

PostgreSQL was chosen because it supports the features the design uses directly: complex CHECK constraints, partial indexes, the JSONB type for logs and measurements, and row locking that skips locked rows (\texttt{SKIP LOCKED}) for safe job claiming. For vector storage, pgvector~\cite{pgvector2026} was suggested at first because it only needs one database, but with an approximate index the filter is applied \emph{after} the index scan, so a query may return fewer results than requested (Section~\ref{sec:3.6.3}). FAISS is only an in-memory library, so the application would have to handle storage, filtering and concurrent access itself. Milvus~\cite{wang2021milvus} was chosen because it filters before searching~\cite{milvus2026filtered}, directly meeting the area-based access control requirement; the trade-off is that Milvus needs etcd and an object store, which makes deployment more complex, although in practice all the services in Docker only use about 515~MiB of memory (Section~\ref{sec:8.7}). Frames are stored in MinIO rather than as files on disk to control access per bucket and keep a hash per object. Milvus already needs an object store, so the application shares that MinIO server in a separate bucket.

\subsection{AI Models}
\label{sec:6.1.4}

\paragraph{Person detection.} Two-stage detectors such as Faster R-CNN are too slow on a CPU, so the project chose the YOLO family. Within YOLO11, the model size directly determines the speed on a CPU (Table~\ref{tab:yolo11-sizes}); even with YOLO11n, the pipeline is about seven times slower than real time (Section~\ref{sec:7.2}), so the project chose the smallest size and compensates for accuracy with a low confidence threshold combined with tracking. YOLO11 is licensed under AGPL-3.0, which suits a non-commercial project, but a commercial product would need a commercial licence; YOLOX~\cite{ge2021yolox}, licensed under Apache-2.0, is therefore registered as a fallback.

\begin{table}[H]
\centering
\caption{Comparison of YOLO11 sizes on the COCO validation set~\cite{ultralytics2024yolo11docs}}
\label{tab:yolo11-sizes}
\begin{tabularx}{\textwidth}{|L|>{\raggedright\arraybackslash}p{2.6cm}|>{\raggedright\arraybackslash}p{3.8cm}|>{\raggedright\arraybackslash}p{3.0cm}|}
\hline
\textbf{Model} & \textbf{mAP$_{50\text{--}95}$} & \textbf{CPU time (ms/image)} & \textbf{Parameters (millions)} \\
\hline
YOLO11n & 39.5 & 56.1 & 2.6 \\
\hline
YOLO11s & 47.0 & 90.0 & 9.4 \\
\hline
YOLO11m & 51.5 & 183.2 & 20.1 \\
\hline
\end{tabularx}
\end{table}

\paragraph{Tracking.} Algorithms that add appearance features, such as DeepSORT~\cite{wojke2017deepsort} and StrongSORT~\cite{du2023strongsort}, keep tracks better through long occlusions, but must run an extra model for every person in every frame, a heavy burden on a CPU in crowded scenes. With static cameras, the project chose ByteTrack~\cite{zhang2022bytetrack} as the default and integrated BoT-SORT~\cite{aharon2022botsort} as an alternative with ReID and camera motion compensation disabled.

\paragraph{Encoder.} Using two separate models for image and text queries would require storing two vectors per appearance and running two models during indexing. General image-text models such as CLIP~\cite{radford2021clip}, including the multilingual variants considered to support Vietnamese, are weaker at separating small differences in clothing. The project chose RaSa~\cite{bai2023rasa} on the supervisor's recommendation: a single model serves all three forms of query with one vector per appearance (Section~\ref{sec:3.5.4}). Its limitation is that queries are English only. The evaluation on the project's data is presented in Section~\ref{sec:8.4}. Part~II of this report studies a different choice for text queries, CLIP ViT-H/14 followed by a reasoning layer, outside the application (Chapter~\ref{chapter:r-overview}).
